\documentclass[11pt]{article}
\usepackage{preamble}

\begin{document}

\packetheading{Lesson 5-4 \textperiodcentered\ Period 2 \textperiodcentered\ Student Packet}{Radical Equations With Rational Exponents + DOK-3 Modeling}

\sectionbanner{OBJECTIVES}
{\small
\renewcommand{\arraystretch}{1.25}
\noindent\begin{tabular}{|p{1.8in}|p{4.8in}|}
\hline
\textbf{Math Objective} & Solve equations with rational exponents by converting to radical form and raising both sides to the reciprocal power; apply rational-exponent models to real-world half-life and decay problems. \\ \hline
\textbf{Language Objective} & Translate a real-world half-life scenario into an equation with a rational exponent, compute, and interpret the result in context. \\ \hline
\textbf{Essential Understanding} & Rational exponents are the bridge between algebraic power operations and real-world exponential/radical modeling --- the half-life formula \(y = a \cdot (1/2)^{(t/h)}\) uses fractional powers directly. \\ \hline
\end{tabular}
}

\sectionbanner{TOPIC GOALS / ESSENTIAL QUESTION / MATERIALS}
{\small
\renewcommand{\arraystretch}{1.25}
\noindent\begin{tabular}{|p{1.8in}|p{4.8in}|}
\hline
\textbf{Topic Goals} & Solve rational-exponent equations; model real-world decay with fractional-power formulas. \\ \hline
\textbf{Essential Question} & How do rational-exponent equations model real-world decay, and when do extraneous solutions arise? \\ \hline
\textbf{Materials} & Student packet, pencil, calculator for Q\#41. Teacher models Ex 4; DOK-3 Practice \#41 is the topic's capstone modeling task. \\ \hline
\end{tabular}
}

\sectionbanner{LAUNCH --- Rational-Exponent Equations (teacher models Example 4)}

\begin{callout}{\IconBook\ RATIONAL-EXPONENT RULES (keep printed on your page)}{calloutgreen}
\textbf{RATIONAL-EXPONENT RULES}
\begin{enumerate}[leftmargin=2.1em,itemsep=1pt,topsep=3pt]
  \item Rewrite \(x^{(m/n)}\) as \((\sqrt[n]{x})^m\) --- bridge to radicals you already know.
  \item To solve \((\text{expr})^{(m/n)} = k\), raise both sides to the reciprocal power \(n/m\).
  \item Simplify, then CHECK --- even rational-exponent equations can have extraneous solutions when the denominator of the exponent is EVEN.
  \item For real-world decay \(y = a \cdot (0.5)^{(t/h)}\): half-life \(h\) tells you the exponent's denominator; plug in \(t\), evaluate.
\end{enumerate}
\end{callout}

\begin{callout}{\IconWarn\ EXTRANEOUS SOLUTIONS ARISE WHEN THE DENOMINATOR IS EVEN}{calloutred}
\((\text{expr})^{(m/n)} = k \rightarrow\) raise both sides to \(n/m\)\\
If \(n\) is EVEN, squaring can introduce solutions that don't satisfy the original equation.\\
ALWAYS substitute back into the ORIGINAL equation to verify.
\end{callout}

\sectionbanner{EXPLORE --- Think-Pair-Share (33 min)}
\textit{Work with a partner. Move in order. Practice \#41 (Half-Life Model) is the big one --- plan \(\sim\)12 min for it.}

\textbf{Bank-item labels (in order):}
\begin{enumerate}[leftmargin=1.6em,itemsep=2pt,topsep=2pt]
  \item Try It 4 --- Rational-exponent equation
  \item Practice \#12 --- Solve with rational exponent
  \item Practice \#33 --- Rational-exponent equation
  \item Practice \#34 --- Rational-exponent equation
  \item Practice \#41  \IconStar\ DOK-3 HALF-LIFE MODEL
\end{enumerate}

\bankitem{Try It 4 --- Rational-exponent equation}{%
Solve each equation.

\begin{enumerate}[label=\textbf{\alph*.},leftmargin=1.8em,itemsep=3pt,topsep=4pt]
  \item \((x^2 - 3x - 6)^{3/2} - 14 = -6\)
  \item \((x - 8)^{3/2} = (10 - x)^3\)
\end{enumerate}
}{}

\bankitem{Practice \#12 --- Solve with rational exponent}{%
\((3x + 2)^{2/3} = 4\)
}{}

\bankitem{Practice \#33 --- Rational-exponent equation}{%
Solve. \((0.5(x^2+5x+136)^{(2)/(3)}=50)\)
}{}

\bankitem{Practice \#34 --- Rational-exponent equation}{%
Solve. \((2(x^2-12x-4)^{(1)/(2)}-3=15)\)
}{}

\bankitem{Practice \#41  \IconStar\ DOK-3 HALF-LIFE MODEL}{%
\textbf{Make Sense and Persevere} The half-life of a certain type of soft drink is 5 h. If you drink 50 mL of this drink, the formula \(y=50(0.5)^{(t)/(5)}\) tells the amount of the drink left in your system after \(t\) hours. How much of the soft drink will be left in your system after 16 hours?
}{}

\sectionbanner{SHARE / SUMMARY (5 min)}

\begin{sentenceframebox}
\textbf{Sentence frames}

Pair share (Half-Life Model): ``This equation has exponent \_\_\_/\_\_\_, so I raise both sides to the power \_\_\_/\_\_\_. After simplifying I get \_\_\_. In the half-life formula, \(a =\) \_\_\_ (start), \(h =\) \_\_\_ (half-life period), \(t =\) \_\_\_ (elapsed time), giving \(y \approx\) \_\_\_.''

Whole class: one pair shares their computation and interpretation. Class signals thumbs up/sideways.
\end{sentenceframebox}

\begin{callout}{\IconSearch\ BRIDGE TO P3}{calloutpeach}
Next period we apply rational-exponent skills to assessment-critical items. Bring your check-for-extraneous habit.
\end{callout}

\sectionbanner{EXIT}

\begin{summaryexitbox}{EXIT TICKET --- Summary Recap}
To solve \((\text{expr})^{(m/n)} = k\) I raise both sides to power \_\_\_.

Half-life tells me the \_\_\_ of the exponent.

After 16 hours, 50 mL of soft drink decays to \_\_\_ mL.

One biggest thing I learned today: \_\_\_\_\_\_\_\_\_\_\_\_\_\_\_\_\_\_\_\_\_\_\_\_\_\_
\end{summaryexitbox}

\clearpage

\sectionbanner{OPTIONAL REINFORCEMENT (not collected)}
\textit{If you finish Explore early. \#19 is a conceptual extraneous-solution argument that extends the half-life skill.}

\bankitem{Optional \#1  (Savvas Practice, extraneous-solution argument)}{%
\textbf{Higher Order Thinking} Some, but not all, equations with rational exponents have extraneous solutions. What is the relationship between the exponents and the possibility of having extraneous solutions for equations with rational exponents? Explain your reasoning.
}{}

\end{document}
